\documentclass[10pt,a4paper]{amsart}
\usepackage[margin=19mm]{geometry}
\usepackage{amsmath,amssymb,mathtools}
\usepackage[scr=rsfs,cal=euler]{mathalfa}
\usepackage{xfrac}
\usepackage[unicode,hidelinks,bookmarks=false]{hyperref}
\newcommand{\ie}{i.e.\ }
\newcommand{\cf}{cf.\ }
\newcommand{\resp}{respectively\ }
\newcommand{\Subsection}{\subsection}
\providecommand{\nobreakdash}{\nobreak\hbox{-}}
\setlength{\emergencystretch}{2em}
\multlinegap0pt
\usepackage{amssymb}
\usepackage{enumerate}
\newtheorem{proposition}{Proposition}
\newcommand{\GL}{\operatorname{GL}}
\newcommand{\Of}{\mathcal{O}_{K,f}}
\newcommand{\Z}{\mathbb{Z}}
\newcommand{\lcm}{\operatorname{lcm}}
\title{The image of the adelic Galois representation of an elliptic curve with complex multiplication}
\author{\'Alvaro Lozano-Robledo, Benjamin York}
\date{Source: arXiv:2603.08545v2 (CC BY 4.0)}
\begin{document}
\maketitle

\noindent\emph{Source and licence.} The image of the adelic Galois representation of an elliptic curve with complex multiplication. Version arXiv:2603.08545v2 (CC BY 4.0), distributed by arXiv under the Creative Commons Attribution 4.0 International licence, http://creativecommons.org/licenses/by/4.0/, as recorded in the arXiv licence metadata for this exact version and shown on the abstract page https://arxiv.org/abs/2603.08545v2. Full licence notice: \texttt{licenses/arxiv-2603.08545-CC-BY-4.0.txt}.\par\smallskip

\noindent\textit{Editorial scope.} Counted unit: the article's proposition prop-normalizerCRT (source0.tex lines 516--518). The article's complete original proof of it is delivered with it (source0.tex lines 519--538, one \texttt{proof} environment of 20 lines). No mathematics is written, repaired or re-derived: the statement and the proof are the article's own lines, in the article's own notation, and no retained line is substituted or re-flowed.  Retained uncounted as the source-local context the proof needs: lines 237--237.  Nothing was omitted from any retained span, so the package carries no omission record.  No retained line contains a citation, so the wrapper carries no bibliography and no bibliography entry.  No retained line contains a figure environment, an image-inclusion command or drawing code, so no image is delivered and the package declares no asset dependency.  Editorial formatting only, in addition: an AMS \texttt{amsart} preamble that restates the article's own declarations and macros used by the retained lines, namely the environments \texttt{proposition} on separate counters, together with the standard packages those lines need.  The retained environments print in this document as Proposition 1 in order of appearance, whereas the article prints the same items with the numbering of its own full document.  The complete native source of the article as distributed in the arXiv e-print for this exact version is bundled unchanged as \texttt{sources/196014/source0.tex}.
\section{Notation and previous results}\label{sec-notation}

\begin{proposition}\label{prop-normalizerCRT}
    Let $\Delta(\Of)=\Delta_K f^2$ be the discriminant of $\Of$ and define the constants $\delta$ and $\phi$ as in Section \ref{sec-notation}. For integers $M, N \geq 2$, let $D = \lcm(M,N)$ and $d = \gcd(M,N)$. Choose $g_{M} \in \mathcal{C}_{\delta, \phi}(M)$ and $g_{N} \in \mathcal{C}_{\delta, \phi}(N)$ so that there is some $g_d \in \mathcal{C}_{\delta, \phi}(d)$ such that $\pi_{M,d}(g_M) = g_d = \pi_{N,d}(g_N)$.    Then, there exists a unique $g_D \in \mathcal{C}_{\delta, \phi}(D)$ such that $\pi_{D,M}(g_D) = g_M$ and $\pi_{D,N}(g_D) = g_N$.
\end{proposition}

\begin{proof} First, we note that if $u,v$ are integers such that $u\equiv v \bmod d$, then the system of congruences 
	\begin{align}\label{eq-1} \begin{cases}
		x\equiv u \bmod M,\\
		x\equiv v \bmod N,
	\end{cases}\end{align} 
	has a unique solution modulo $x_0=x_0(u,v)\bmod D$. Indeed, if we write $M=M'd$ and $N=N'd$, then the system (\ref{eq-1}) is equivalent to 
		\begin{align}\label{eq-2} \begin{cases}
			x\equiv u \bmod M',\\
			x\equiv u \bmod d,\\
			x\equiv v \bmod N',\\
			x\equiv v \bmod d
	\end{cases}\end{align}
which, by the Chinese remainder theorem, has a unique solution modulo $M'N'd = D$ if and only if $u\equiv v \bmod d$. 

In particular, the uniqueness of solutions of the system (\ref{eq-1}) modulo $D$ shows that for any $g_N \in \GL(2, \Z/N\Z)$ and $g_M\in \GL(2, \Z/M\Z)$, there exists a unique $g_D \in \GL(2,\Z/D\Z)$ which reduces to $g_N$ and $g_M$ modulo $N$ and $M$, respectively, if and only if $g_N\equiv g_M \bmod d$.

Now, if  $g_{M} \in \mathcal{C}_{\delta, \phi}(M)$ and $g_{N} \in \mathcal{C}_{\delta, \phi}(N)$ so that there is some $g_d \in \mathcal{C}_{\delta, \phi}(d)$ such that $\pi_{M,d}(g_M) = g_d = \pi_{N,d}(g_N)$, as in the statement of the lemma, then our previous remarks show that there is a unique $g_D\in \GL(2,\Z/D\Z)$ such that $g_D\equiv g_M\bmod M$ and $g_D\equiv g_N\bmod N$. Moreover, since $g_M$ is an element of the Cartan subgroup, it must be of the form $c_{\delta,\phi}(a_M,b_M)$ for some integers $a_M,b_M$ (as in the notation established in Section \ref{sec-notation}), and similarly $g_N= c_{\delta,\phi}(a_N,b_N)$ for some $a_N,b_N$. Let $x_0(a_M,a_N)$ and $x_0(b_M,b_N)$ be the unique solutions of the corresponding systems as in Eq. (\ref{eq-1}).  Then, the element
$$g=c_{\delta,\phi}(x_0(a_M,a_N),x_0(b_M,b_N))\in  \mathcal{C}_{\delta, \phi}(D)\subseteq \GL(2,\Z/D\Z)$$
satisfies $g\equiv g_N \bmod N$ and $g\equiv g_M\bmod M$, and therefore, by the uniqueness of $g_D$, we must have $g=g_D$ belongs to $\mathcal{C}_{\delta, \phi}(D)$, as desired. 
 \end{proof}

\end{document}
